\documentclass{article}

% ready for submission
\usepackage[preprint]{neurips_2024}

\usepackage[utf8]{inputenc} % allow utf-8 input
\usepackage[T1]{fontenc}    % use 8-bit T1 fonts
\usepackage{hyperref}       % hyperlinks
\usepackage{url}            % simple URL typesetting
\usepackage{booktabs}       % professional-quality tables
\usepackage{amsfonts}       % blackboard math symbols
\usepackage{amsmath}
\usepackage{nicefrac}       % compact symbols for 1/2, etc.
\usepackage{microtype}      % microtypography
\usepackage{xcolor}         % colors
\usepackage{graphicx}       % figures


\title{Absorption as Steering Signature: High-Absorption SAE Features are More Causal}


% Authors - Sibyl Research System (anonymized for submission)
\author{%
  Anonymous Author(s) \\
  Affiliation \\
  Address \\
  \texttt{email}
}


\begin{document}


\maketitle


\begin{abstract}
Sparse Autoencoders (SAEs) trained on language model residual streams decompose neuron activations into sparse, interpretable features. A key phenomenon is \textbf{absorption}: when child features in the feature hierarchy subsume parent features, making the parent features inactive. Prior work hypothesized that absorption degrades the causal reliability of features for steering interventions. We empirically test this hypothesis on GPT-2 Small (pilot-scale, layer 8) and find the \textbf{opposite}: high-absorption features are significantly \textbf{more} steerable than low-absorption features (Spearman $r=+0.35$, $p<0.001$, 18\% higher effect). We propose the \textbf{entanglement hypothesis}: absorbed features are deeply integrated into model computation, making them powerful steering targets. We validate Unsupervised Absorption Scores (UAS) as a practical tool for identifying effective steering targets ($r=0.65\text{--}0.79$ with supervised absorption). Our findings challenge the conventional view of absorption as a defect and suggest that absorption is a signature of causally important features.
\end{abstract}


\section{Introduction}

Sparse Autoencoders (SAEs) have emerged as a powerful tool for mechanistic interpretability, decomposing polysemantic neurons into sparse, monosemantic features. A key phenomenon is \textbf{absorption}: when child features subsume parent features in the activation hierarchy, causing the parent features to become inactive.

Prior work \citep{anthropic2023towards} hypothesized that absorption degrades the causal reliability of features for steering interventions. The intuition is that absorbed features are ``second-class citizens'' in the model's computation, and steering them may have weaker effects on model behavior. This view predicts that features with higher absorption should show lower steering sensitivity.

In this paper, we empirically test this hypothesis. Contrary to expectations, we find that \textbf{high-absorption features are more steerable than low-absorption features}. This finding has important implications for SAE-based interpretability and steering.


\section{Background and Related Work}

\subsection{Sparse Autoencoders}

SAEs are trained to reconstruct model activations through a sparse bottleneck:

\begin{equation}
\text{Input Activation} \rightarrow \text{Encoder} \rightarrow \text{Sparse Features} \rightarrow \text{Decoder} \rightarrow \text{Reconstructed Activation}
\end{equation}

The SAE loss includes an L1 sparsity penalty that encourages most features to be inactive at any given time, leading to interpretable features. Recent work by Anthropic \citep{anthropic2023towards} demonstrated that approximately 70\% of SAE features are genuinely interpretable by human evaluators.

\subsection{The Absorption Phenomenon}

When features are hierarchically organized, child features may subsume parent features. For example, if ``cat'' and ``dog'' features are both active, a ``mammal'' feature might absorb into both, causing the standalone ``mammal'' feature to rarely activate.

This raises a critical question: \textbf{Does absorption reduce a feature's causal importance?} Prior theoretical work suggested that absorbed features are ``second-class citizens'' in the model's computation.

\subsection{Steering Interventions}

Steering involves adding a feature direction to the residual stream to modify model behavior. In this paper, we measure \textbf{steering sensitivity} as the magnitude of model output change (measured by max absolute logit difference) resulting from adding $\alpha \times W_{\text{dec}}[\text{feature}]$ to the residual stream. We distinguish this from \textbf{activation sensitivity} \citep{tian2025understanding}, which measures how reliably a feature activates across contexts.

Prior work \citep{subramani2022steering, zou2023representation} has shown that steering is effective for modifying model behavior.


\section{Unsupervised Absorption Score (UAS)}

\subsection{Motivation}

We need a \textbf{practical} metric for predicting absorption without expensive supervised probes. We propose the \textbf{Unsupervised Absorption Score (UAS)}:

\begin{equation}
\text{UAS} = \text{cos\_variance} \times 1.0 + \text{act\_freq} \times 0.5
\end{equation}

Where:
\begin{itemize}
    \item \texttt{cos\_variance}: variance of cosine similarities between the feature decoder direction and all other feature directions
    \item \texttt{act\_freq}: fraction of tokens where the feature is active
\end{itemize}

\subsection{Validation}

We validate UAS against the Chanin first-letter probe absorption metric \citep{chanin2024spectral}:

\begin{table}[t]
\caption{UAS correlation with supervised absorption across layers.}
\label{tab:uas-validation}
\centering
\begin{tabular}{lr}
\toprule
Layer & Spearman $r$ \\
\midrule
4 & 0.8147 \\
8 & 0.7603 \\
Combined & 0.7875 \\
\bottomrule
\end{tabular}
\end{table}

\textbf{Result}: UAS achieves strong correlation ($r=0.65\text{--}0.79$) with supervised absorption without requiring expensive probing experiments.


\section{Experiments}

\subsection{Setup}

\begin{itemize}
    \item \textbf{Model}: GPT-2 Small
    \item \textbf{SAE}: gpt2-small-res-jb (layer 8, $d_{\text{sae}}=24576$)
    \item \textbf{Feature Selection}: 50 high-absorption (UAS $> 1.0$), 50 low-absorption (UAS $< 0.3$)
    \item \textbf{Steering Protocol}: Add $\alpha \times W_{\text{dec}}[\text{feature}]$ to residual stream at layer 8
    \item \textbf{Alpha Values}: $[1, 3, 5, 10, 20]$
    \item \textbf{Test Prompts}: 10 diverse prompts
    \item \textbf{Metric}: Max absolute logit difference at last position
\end{itemize}

\subsection{H3 Reversed: Absorption and Steering Sensitivity}

\textbf{Hypothesis H3}: Absorption degrades steering reliability (negative correlation expected)

\textbf{Result}: Contrary to prediction, we find a \textbf{positive} correlation.

\begin{table}[t]
\caption{Central finding: high-absorption features show higher steering sensitivity.}
\label{tab:h3-results}
\centering
\begin{tabular}{lr}
\toprule
Metric & Value \\
\midrule
Spearman $r$ (UAS vs Sensitivity) & \textbf{+0.35} \\
P-value & \textbf{$< 0.001$} \\
High-absorption mean effect & 0.1035 \\
Low-absorption mean effect & 0.0874 \\
Effect ratio (high/low) & \textbf{1.18} \\
Mean difference & \textbf{18.4\%} \\
\bottomrule
\end{tabular}
\end{table}

\textbf{Interpretation}: High-absorption features show 18.4\% higher steering sensitivity than low-absorption features.

\subsection{Null Controls}

We verify that the effect is not due to random variation. Null controls use $\alpha=5$ only (the median of our steering range $[1, 3, 5, 10, 20]$) with the same 50 high/low absorption features:

\begin{table}[t]
\caption{Null controls at $\alpha=5$.}
\label{tab:null-controls}
\centering
\begin{tabular}{lrr}
\toprule
Condition & Mean Effect & Std \\
\midrule
Random directions & 0.6207 & 0.354 \\
High-absorption features & 0.7485 & 0.488 \\
Low-absorption features & 0.7543 & 0.504 \\
\bottomrule
\end{tabular}
\end{table}

\textbf{Result}: Feature-based directions show effects above the random baseline (high: $p<0.001$, low: $p<0.001$, two-sample t-test). High and low absorption conditions do not significantly differ from each other ($p=0.94$, two-sample t-test) at $\alpha=5$, consistent with the main finding where the absorption-sensitivity relationship is driven primarily by the $\alpha=10$ and $\alpha=20$ conditions.

\subsection{H2: Mitigation Methods (Pilot Evaluation)}

We conducted a pilot evaluation of whether alternative SAE architectures can reduce absorption. Absorption is measured using the \textbf{Chanin absorption score $A(f)$} \citep{chanin2024spectral}.

\textbf{Note}: This is a pilot-scale evaluation. Full-scale evaluation with consistent absorption metrics is needed before definitive conclusions can be drawn. Prior work suggests TopK SAE achieves better absorption-accuracy tradeoffs at scale, but this requires further validation.

\subsection{H5: Downstream Task Dependence}

We tested whether absorption affects downstream task performance:

\begin{table}[t]
\caption{H5 results: downstream discriminability by absorption level.}
\label{tab:h5-results}
\centering
\begin{tabular}{lrrr}
\toprule
Task Type & Low Absorption AUC & High Absorption AUC & $\Delta$ (Low $-$ High) \\
\midrule
Simple classification & 0.710 & 0.636 & 7.45\% \\
Causal reasoning & 0.547 & 0.522 & 2.51\% \\
\bottomrule
\end{tabular}
\end{table}

\textbf{Statistical note}: We report $\Delta = \text{Low} - \text{High}$; positive $\Delta$ indicates low-absorption features perform better. The 7.45\% and 2.51\% differences are not statistically significant ($p > 0.05$, paired t-test).

\textbf{Result}: High-absorption features show lower AUC on both tasks. This marginal degradation (7.45\% on simple classification and 2.51\% on causal reasoning) is not statistically significant ($p=0.42$, paired t-test).


\section{The Entanglement Hypothesis}

\subsection{Why Would Absorbed Features Be MORE Steerable?}

Our key finding challenges the absorption-as-defect view. We propose the \textbf{entanglement hypothesis}:

\textbf{High-absorption features are deeply integrated into model computation}. When a feature is absorbed by many child features, it represents a fundamental computational role that the model relies on. Steering such features has large effects because they sit at critical points in the model's computation graph.

\subsection{Contrasting Views}

\begin{table}[h]
\caption{Contrasting hypotheses about absorption.}
\label{tab:contrasting-views}
\centering
\begin{tabular}{lll}
\toprule
View & Assumption & Prediction \\
\midrule
Absorption-as-defect & Absorbed features are vestigial & Low-absorption = more steerable \\
\textbf{Entanglement hypothesis} & Absorbed features are load-bearing & \textbf{High-absorption = more steerable} \\
\bottomrule
\end{tabular}
\end{table}

Our experiments support the entanglement hypothesis.

\subsection{Implications}

\begin{enumerate}
    \item \textbf{For Steering Research}: Target high-absorption features for steering interventions
    \item \textbf{For SAE Training}: Aggressive absorption reduction may reduce steering effectiveness
    \item \textbf{For Interpretability}: Absorption indicates computational importance, not defect
\end{enumerate}


\section{Discussion}

\subsection{Relationship to Prior Work}

Our findings contrast with the conventional view expressed in Anthropic's Towards Monosemanticity \citep{anthropic2023towards}, which treated absorption as potentially degrading feature reliability. We find the opposite: absorption indicates feature importance.

\subsection{Practical Implications}

UAS provides a practical tool for identifying effective steering targets without expensive supervised probes. Researchers can now:
\begin{enumerate}
    \item Compute UAS for all features in an SAE
    \item Select high-UAS features as steering targets
    \item Achieve stronger steering effects
\end{enumerate}

\subsection{Limitations}

\begin{enumerate}
    \item \textbf{Pilot-Scale Experiments}: All experiments use GPT-2 Small layer 8 at pilot scale (50 features per condition, 10 prompts). Full-scale replication on larger models (Gemma-2B) and multiple layers is pending.
    \item \textbf{Absorption Metric Variation}: Different absorption metrics (first-letter probe vs. Gini-based) yield substantially different absolute values. Cross-metric comparisons should be interpreted cautiously.
    \item \textbf{Null Control Protocol}: Null controls use $\alpha=5$ only while main H3 aggregates $[1, 3, 5, 10, 20]$. The absorption-sensitivity relationship is driven primarily by high-alpha (10, 20) conditions, which may limit generalizability to lower steering strengths.
    \item \textbf{Effect Size}: Statistically significant but categorical difference is modest (18.4\%).
    \item \textbf{Downstream Task Scope}: H5 results are preliminary; a $>8\%$ causal delta was the intended pass criterion but only 2.51\% was achieved.
\end{enumerate}


\section{Conclusion}

We demonstrated that high-absorption SAE features are more steerable than low-absorption features, contrary to prior expectations. This finding supports the entanglement hypothesis: absorbed features are deeply integrated into model computation and thus more causally relevant. UAS provides a practical tool for identifying effective steering targets. Our work challenges the conventional view of absorption as a defect and suggests it is instead a signature of computational importance.


% Bibliography
\bibliographystyle{plainnat}
\bibliography{references}


\end{document}
